\begin{scanpage}{0029}{12}
\sectiontitle{2}{Asymptotic Expressions}
The integral
\[v(r,\psi)=\frac{i}{2\alpha}\sin\frac{\pi^2}{\alpha}\int_{D_0}\frac{e^{ikr\cos\zeta}\,d\zeta}{\cos\frac{\pi^2}{\alpha}-\cos\frac\pi\alpha(\psi+\zeta)},\tag{2.01}\]
in Equations (1.31) and (1.32) generally is not expressed in terms of well-known functions. However, when $kr\gg1$ it may be calculated approximately by the method of steepest descents [21]. In Integral (2.01), changing for this purpose to a new integration variable
\[s=\sqrt2e^{i\frac\pi4}\sin\frac\zeta2,\quad s^2=i(1-\cos\zeta),\]
we obtain
\[v(r,\psi)=\frac{\sin\frac\pi n}{\sqrt2\pi n}e^{i(kr+\frac\pi4)}\int_{-\infty}^{\infty}\frac{e^{-krs^2}\,ds}{\left(\cos\frac\pi n-\cos\frac{\psi+\zeta}{n}\right)\cos\frac\zeta2},\tag{2.02}\]
where
\[n=\frac\alpha\pi.\tag{2.03}\]
It is not difficult to see that the point $s=0$ is a saddle point: as one goes further from it along the imaginary axis $(\operatorname{Re}s=0)$ in the plane of the complex variable $s$, the function $e^{-krs^2}$ most rapidly increases, and as one goes along the real axis $(\operatorname{Im}s=0)$ it decreases most rapidly. Therefore, when $kr\gg1$ the main contribution to the integral (2.02) is given by the integrand in the section of the contour in the vicinity of the saddle point $(s=0)$.

The method of steepest descents is carried out by expanding the integrand (except for the factor $e^{-krs^2}$) into a Taylor series in powers of $s$. This series is then integrated term by term. If the integrand
\end{scanpage}
